% Chapter 6 — ZK Password-Locked UTxO from Scratch.
% All source code lives in zk-password/ at the root of the repository.

% ─────────────────────────────────────────────────────────────────────────────
\section{What We Are Building} \label{sec:ch6-overview}

This chapter builds the simplest possible ZK smart contract on Cardano: a UTxO
that can only be unlocked by someone who can \emph{prove} they know a secret
password --- without ever revealing the password on-chain.

The password is hashed with \textbf{Poseidon}, a ZK-friendly hash function. The
hash is stored publicly in the datum. To spend the UTxO, the spender generates a
Groth16 proof that they know a pre-image of that hash. The Aiken validator
verifies the proof on-chain using the BLS12-381 pairing primitives introduced
in Plutus V3.

\begin{highlight}
\textbf{What stays private:} the password. \quad
\textbf{What goes on-chain:} the Poseidon hash, the verification key, and the proof.
\end{highlight}

The complete source lives in \texttt{zk-password/} at the root of the repository.
By the end of this chapter the directory will look like this:

\begin{verbatim}
zk-password/
+-- aiken.toml
+-- package.json
+-- circuit/
|   +-- circuit.circom          <- Circom constraint circuit
|   +-- poseidon_hash.circom    <- helper: compute hash without constraint
|   +-- gen_input.js            <- script: password -> input.json
|   +-- compress_vk.js          <- script: compress verification key
|   +-- compress_proof.js       <- script: compress proof points
+-- lib/
|   +-- zk_password/
|       +-- groth16.ak          <- Groth16 verifier library
|       +-- password_tests.ak   <- Aiken unit tests
+-- validators/
    +-- password.ak             <- on-chain validator
\end{verbatim}

% ─────────────────────────────────────────────────────────────────────────────
\section{Project Setup} \label{sec:ch6-setup}

Create the Aiken project and initialise npm for the JavaScript utility scripts:

\begin{verbatim}
aiken new chapter6/zk-password
cd zk-password
npm init -y
npm install -D circomlib
npm install @noble/curves
\end{verbatim}

\texttt{circomlib} provides the Poseidon circuit templates.
\texttt{@noble/curves} is used by the compression scripts to convert uncompressed
BLS12-381 points to the 48-byte (G1) and 96-byte (G2) format that Cardano accepts.

% ─────────────────────────────────────────────────────────────────────────────
\section{The Circuit} \label{sec:ch6-circuit}

Create \texttt{circuit/circuit.circom}:

\begin{verbatim}
pragma circom 2.1.6;

include "../node_modules/circomlib/circuits/poseidon.circom";

// Prove knowledge of a password whose Poseidon hash matches a commitment.
//
// Public inputs  (visible on-chain in the datum):
//   hash  -- Poseidon(pwd)
//
// Private inputs (never leave the prover's device):
//   pwd   -- the password
template PasswordLock() {
    signal input hash;   // public
    signal input pwd;    // private

    component h = Poseidon(1);
    h.inputs[0] <== pwd;

    hash === h.out;
}

component main {public [hash]} = PasswordLock();
\end{verbatim}

The \texttt{===} operator defines an arithmetic constraint: if the Poseidon hash
of \texttt{pwd} does not equal \texttt{hash}, the constraint system is
unsatisfiable and no valid proof can ever be constructed.
The \texttt{<==} operator assigns a value and simultaneously constrains the signal.

\subsection*{Compile the Circuit}

From inside \texttt{zk-password/}:

\begin{verbatim}
circom --r1cs --wasm --sym --inspect \
       --prime bls12381 circuit/circuit.circom -o circuit/
\end{verbatim}

The \texttt{--prime bls12381} flag targets the BLS12-381 scalar field, which
matches Cardano's on-chain pairing primitives. The outputs land in \texttt{circuit/}:

\begin{itemize}
  \item \texttt{circuit.r1cs} --- the binary constraint system.
  \item \texttt{circuit\_js/circuit.wasm} --- the WebAssembly witness calculator.
  \item \texttt{circuit.sym} --- the symbol table (signal names and witness indices).
\end{itemize}

\begin{figure}[h]
  \centering
  \includegraphics[width=0.85\linewidth]{circom_compiled.png}
  \caption{Circom compilation output. Our circuit has 216 non-linear constraints,
           1 public input (\texttt{hash}), and 1 private input (\texttt{pwd}).}
\end{figure}

Inspect the compiled R1CS:

\begin{verbatim}
snarkjs r1cs info circuit/circuit.r1cs
\end{verbatim}

You should see \texttt{Curve: bls12-381}, 1 public input, 1 private input, and
roughly 240 constraints (mostly Poseidon round internals).

% ─────────────────────────────────────────────────────────────────────────────
\section{Generating \texttt{input.json}} \label{sec:ch6-input}

The circuit needs two inputs: the password as a field element and its Poseidon
hash. There is a subtlety here worth understanding before writing any code.

\subsection*{The Field-Element Problem}

The \texttt{pwd} signal in the circuit is a BLS12-381 field element --- an
integer, not a string. A plain numeric string like \texttt{"12345"} converts
trivially with \texttt{BigInt("12345")}. A word like \texttt{"mypassword"}
does not, and calling \texttt{BigInt("mypassword")} throws a runtime error.

The solution is to encode the string as a big-endian integer of its UTF-8 bytes:

\begin{verbatim}
function toFieldElement(s) {
  if (/^\d+$/.test(s)) return BigInt(s);
  return BigInt("0x" + Buffer.from(s, "utf8").toString("hex"));
}
\end{verbatim}

\subsection*{The Hash Computation Problem}

We also need the Poseidon hash of the password, but we cannot just call
\texttt{circomlibjs}'s \texttt{buildPoseidon} here. That library implements
Poseidon over the \textbf{BN128} field; our circuit uses the \textbf{BLS12-381}
scalar field. The two fields have different primes, so the hash values differ.

We solve this by compiling a tiny helper circuit that simply exposes the hash
as an output signal --- no equality constraint, just a wire:

\textbf{\texttt{circuit/poseidon\_hash.circom}}:

\begin{verbatim}
pragma circom 2.1.6;

include "../node_modules/circomlib/circuits/poseidon.circom";

template HashOnly() {
    signal input  pwd;
    signal output hash;
    component h = Poseidon(1);
    h.inputs[0] <== pwd;
    hash <== h.out;
}

component main = HashOnly();
\end{verbatim}

Compile it with the same BLS12-381 prime:

\begin{verbatim}
circom --wasm --prime bls12381 \
       circuit/poseidon_hash.circom -o circuit/
\end{verbatim}

Because this circuit has no \texttt{===} constraint, the witness calculator runs
cleanly for any input and \texttt{witness[1]} --- the first public output ---
contains \texttt{Poseidon(pwd)} computed over the correct field.

\subsection*{\texttt{circuit/gen\_input.js}}

\begin{verbatim}
const { readFileSync, writeFileSync } = require("fs");
const path = require("path");

const password = process.argv[2] ?? "12345";

function toFieldElement(s) {
  if (/^\d+$/.test(s)) return BigInt(s);
  return BigInt("0x" + Buffer.from(s, "utf8").toString("hex"));
}

const pwd = toFieldElement(password);

async function main() {
  const buildWC = require("./poseidon_hash_js/witness_calculator.js");
  const wasm = readFileSync(
    path.join(__dirname, "poseidon_hash_js", "poseidon_hash.wasm")
  );
  const wc = await buildWC(wasm);

  const witness = await wc.calculateWitness({ pwd }, true);
  const hash = witness[1]; // output signal is always at index 1

  const input = { hash: hash.toString(), pwd: pwd.toString() };
  const outPath = path.join(__dirname, "input.json");
  writeFileSync(outPath, JSON.stringify(input, null, 2));

  console.log("password :", password);
  console.log("hash     :", hash.toString());
  console.log("written  :", outPath);
}

main().catch(err => { console.error(err); process.exit(1); });
\end{verbatim}

Usage:

\begin{verbatim}
node circuit/gen_input.js 123456
node circuit/gen_input.js mypassword
node circuit/gen_input.js "hello world"
\end{verbatim}

The resulting \texttt{circuit/input.json} contains the hash and password as
quoted decimal strings, ready for \texttt{snarkjs groth16 fullprove}.

% ─────────────────────────────────────────────────────────────────────────────
\section{The Trusted Setup} \label{sec:ch6-trusted-setup}

Groth16 requires a two-phase trusted setup. Phase 1 (Powers of Tau) is
circuit-independent; Phase 2 is specific to your R1CS.

\subsection*{Phase 1 --- Powers of Tau}

\begin{verbatim}
snarkjs powersoftau new bls12381 12 pot12_0000.ptau -v
snarkjs powersoftau contribute pot12_0000.ptau pot12_0001.ptau \
    --name="dev" -e="some entropy"
snarkjs powersoftau beacon pot12_0001.ptau pot12_beacon.ptau \
    00000000000000000000000000000000 10 -n="Dev beacon"
snarkjs powersoftau prepare phase2 pot12_beacon.ptau pot12_final.ptau -v
\end{verbatim}

We use $2^{12} = 4096$ max constraints --- well above our 216-constraint circuit
and fast to generate. Use a larger power (e.g.\ 14) for bigger circuits.

\subsection*{Phase 2 --- Circuit-Specific Setup}

\begin{verbatim}
snarkjs groth16 setup circuit/circuit.r1cs pot12_final.ptau zkey_0000.zkey
snarkjs zkey contribute zkey_0000.zkey zkey_0001.zkey \
    --name="dev" -e="more entropy"
snarkjs zkey beacon zkey_0001.zkey zkey_final.zkey \
    00000000000000000000000000000000 10 -n="Final beacon"
snarkjs zkey export verificationkey zkey_final.zkey verification_key.json
\end{verbatim}

\begin{figure}[h]
  \centering
  \includegraphics[width=0.85\linewidth]{verification_key_generation.png}
  \caption{snarkjs exports the verification key to \texttt{verification\_key.json}.
           This file contains the uncompressed BLS12-381 G1 and G2 points.}
\end{figure}

\begin{remark}
For production you should use an existing BLS12-381 Powers of Tau file
(e.g.\ from the Zcash ceremony) and run at least three independent Phase 2
contributions. The parameters here are safe only for local development.
\end{remark}

% ─────────────────────────────────────────────────────────────────────────────
\section{Generating the Proof} \label{sec:ch6-prove}

With \texttt{circuit/input.json}, the compiled WASM, and the final zkey, generate
the proof in one command:

\begin{verbatim}
snarkjs groth16 fullprove \
    circuit/input.json \
    circuit/circuit_js/circuit.wasm \
    zkey_final.zkey \
    proof.json public.json
\end{verbatim}

Verify the proof off-chain before going further:

\begin{verbatim}
snarkjs groth16 verify verification_key.json public.json proof.json
\end{verbatim}

You should see \texttt{[INFO] snarkJS: OK!}. The \texttt{public.json} file
contains the single public input: the Poseidon hash of your password.

% ─────────────────────────────────────────────────────────────────────────────
\section{Compressing Points for Cardano} \label{sec:ch6-compress}

Cardano's BLS12-381 builtins only accept \textbf{compressed} elliptic-curve
points. The snarkjs output uses uncompressed Weierstrass coordinates
$[x, y, 1]$ for G1 and $[[x_0, x_1], [y_0, y_1], [1, 0]]$ for G2.
Compressed form encodes a point in its $x$-coordinate alone (48 bytes for G1,
96 bytes for G2), using one flag bit to disambiguate which of the two possible
$y$ values to use.

We compress both the verification key and the proof with the
\texttt{@noble/curves} library.

\subsection*{\texttt{circuit/compress\_vk.js}}

\begin{verbatim}
const { readFileSync, writeFileSync } = require("fs");
const path = require("path");

const vkPath  = process.argv[2]
    ?? path.join(__dirname, "..", "verification_key.json");
const outPath = process.argv[3]
    ?? path.join(__dirname, "..", "compressed_vk.json");

async function main() {
  const { bls12_381 } = await import("@noble/curves/bls12-381.js");
  const G1 = bls12_381.G1.Point;
  const G2 = bls12_381.G2.Point;

  const vk = JSON.parse(readFileSync(vkPath, "utf8"));

  function compressG1([x, y]) {
    return Buffer.from(
      G1.fromAffine({ x: BigInt(x), y: BigInt(y) }).toBytes(true)
    ).toString("hex");
  }

  function compressG2([[x0, x1], [y0, y1]]) {
    return Buffer.from(
      G2.fromAffine({
        x: { c0: BigInt(x0), c1: BigInt(x1) },
        y: { c0: BigInt(y0), c1: BigInt(y1) },
      }).toBytes(true)
    ).toString("hex");
  }

  const compressed = {
    vk_alpha1: compressG1(vk.vk_alpha_1),
    vk_beta2:  compressG2(vk.vk_beta_2),
    vk_gamma2: compressG2(vk.vk_gamma_2),
    vk_delta2: compressG2(vk.vk_delta_2),
    vk_ic:     vk.IC.map(compressG1),
  };

  writeFileSync(outPath, JSON.stringify(compressed, null, 2));
}

main().catch(err => { console.error(err); process.exit(1); });
\end{verbatim}

\begin{verbatim}
node circuit/compress_vk.js
\end{verbatim}

The resulting \texttt{compressed\_vk.json} has:

\begin{itemize}
  \item \texttt{vk\_alpha1} --- 96 hex chars (48 bytes, compressed G1)
  \item \texttt{vk\_beta2}, \texttt{vk\_gamma2}, \texttt{vk\_delta2} ---
        192 hex chars each (96 bytes, compressed G2)
  \item \texttt{vk\_ic} --- array of two 96-hex-char strings; the first element
        is the constant term $IC_0$, the second is $IC_1$ (one entry per public input)
\end{itemize}

\subsection*{\texttt{circuit/compress\_proof.js}}

The proof points use the same approach:

\begin{verbatim}
const { readFileSync, writeFileSync } = require("fs");
const path = require("path");

const proofPath = process.argv[2]
    ?? path.join(__dirname, "..", "proof.json");
const outPath   = process.argv[3]
    ?? path.join(__dirname, "..", "compressed_proof.json");

async function main() {
  const { bls12_381 } = await import("@noble/curves/bls12-381.js");
  const G1 = bls12_381.G1.Point;
  const G2 = bls12_381.G2.Point;

  const proof = JSON.parse(readFileSync(proofPath, "utf8"));

  const compressed = {
    pi_a: Buffer.from(
      G1.fromAffine({
        x: BigInt(proof.pi_a[0]),
        y: BigInt(proof.pi_a[1]),
      }).toBytes(true)
    ).toString("hex"),
    pi_b: Buffer.from(
      G2.fromAffine({
        x: { c0: BigInt(proof.pi_b[0][0]), c1: BigInt(proof.pi_b[0][1]) },
        y: { c0: BigInt(proof.pi_b[1][0]), c1: BigInt(proof.pi_b[1][1]) },
      }).toBytes(true)
    ).toString("hex"),
    pi_c: Buffer.from(
      G1.fromAffine({
        x: BigInt(proof.pi_c[0]),
        y: BigInt(proof.pi_c[1]),
      }).toBytes(true)
    ).toString("hex"),
  };

  writeFileSync(outPath, JSON.stringify(compressed, null, 2));
}

main().catch(err => { console.error(err); process.exit(1); });
\end{verbatim}

\begin{verbatim}
node circuit/compress_proof.js
\end{verbatim}

% ─────────────────────────────────────────────────────────────────────────────
\section{The Groth16 Verifier Library} \label{sec:ch6-groth16}

The file \texttt{lib/zk\_password/groth16.ak} is a self-contained Groth16
verifier modelled after Janus Wallet's implementation.

\subsection*{Types}

\begin{verbatim}
pub type VerificationKey {
  alpha: ByteArray,        // G1 — 48 bytes
  beta:  ByteArray,        // G2 — 96 bytes
  gamma: ByteArray,        // G2 — 96 bytes
  delta: ByteArray,        // G2 — 96 bytes
  ic:    List<ByteArray>,  // G1 points; length = num_public_inputs + 1
}

pub type Proof {
  pi_a: ByteArray,  // G1 — 48 bytes
  pi_b: ByteArray,  // G2 — 96 bytes
  pi_c: ByteArray,  // G1 — 48 bytes
}
\end{verbatim}

All fields are compressed bytes. The \texttt{ic} list has two elements for a
circuit with one public input: the constant term $IC_0$ and the coefficient
$IC_1$ for \texttt{hash}.

\subsection*{The Verification Function}

The pairing equation that Groth16 checks is:

\[
  e(\pi_A,\, \pi_B) = e(\alpha, \beta)\cdot e(\mathbf{vk\_x},\, \gamma)\cdot e(\pi_C,\, \delta)
\]

where $\mathbf{vk\_x} = IC_0 + \sum_i (\text{public\_input}_i \cdot IC_{i+1})$.

\begin{verbatim}
pub fn groth16_verify(
  vk: VerificationKey,
  proof: Proof,
  public_inputs: List<Scalar>,
) -> Bool {
  let a = g1.decompress(proof.pi_a)
  let b = g2.decompress(proof.pi_b)
  let c = g1.decompress(proof.pi_c)

  let vk_x =
    when vk.ic is {
      [] -> fail @"vk.ic must not be empty"
      [ic0, ..rest] -> derive(rest, public_inputs, g1.decompress(ic0))
    }

  let e_ab         = bls12_381_miller_loop(a, b)
  let e_alpha_beta =
    bls12_381_miller_loop(g1.decompress(vk.alpha), g2.decompress(vk.beta))
  let e_x_gamma    = bls12_381_miller_loop(vk_x, g2.decompress(vk.gamma))
  let e_c_delta    = bls12_381_miller_loop(c, g2.decompress(vk.delta))

  let rhs =
    bls12_381_mul_miller_loop_result(
      bls12_381_mul_miller_loop_result(e_alpha_beta, e_x_gamma),
      e_c_delta,
    )

  bls12_381_final_verify(e_ab, rhs)
}
\end{verbatim}

The entire cryptographic work happens inside Cardano's builtin functions
(\texttt{bls12\_381\_miller\_loop} and \texttt{bls12\_381\_final\_verify}).
These are native operations in Plutus V3 and are priced at fixed, low execution
units regardless of proof complexity.

% ─────────────────────────────────────────────────────────────────────────────
\section{The Aiken Validator} \label{sec:ch6-validator}

\subsection*{The Opaque Scalar Problem}

A natural first attempt for the datum would be:

\begin{verbatim}
pub type Datum {
  credential_hash: Scalar,  -- would be natural, but WRONG
  vk: VerificationKey,
}
\end{verbatim}

The Aiken compiler immediately rejects this:

\begin{verbatim}
x I caught an opaque type trying to escape
  |
  | datum: Option<Datum>,
  |        ──────────────  invalid validator's boundary
  |
  = Opaque types cannot figure anywhere in an outward-facing type
    like a validator's redeemer or datum.
\end{verbatim}

\texttt{Scalar} is \emph{opaque}: its constructor is hidden and it enforces
invariants that cannot be expressed as plain Plutus Data. Opaque types cannot
cross the on-chain/off-chain boundary (datum or redeemer), because the
off-chain world has no way to build them safely from raw bytes.

\textbf{The fix:} store the hash as a plain \texttt{Int} in the datum (a
normal Plutus integer that any off-chain tool can construct), and convert it
to a \texttt{Scalar} at runtime inside the validator body using
\texttt{scalar.new}:

\begin{verbatim}
use aiken/crypto/bls12_381/scalar
use cardano/transaction.{OutputReference, Transaction}
use zk_password/groth16.{Proof, VerificationKey, groth16_verify}

pub type Datum {
  credential_hash: Int,   // Poseidon(pwd) as a plain integer
  vk: VerificationKey,    // compressed Groth16 verification key
}

pub type Redeemer {
  proof: Proof,           // compressed Groth16 proof
}

validator password_lock {
  spend(
    datum: Option<Datum>,
    redeemer: Redeemer,
    _self: OutputReference,
    _ctx: Transaction,
  ) {
    expect Some(d) = datum
    // Convert Int -> Scalar at runtime. scalar.new returns None if
    // the value is outside [0, field_prime), which rejects malformed datums.
    expect Some(hash_scalar) = scalar.new(d.credential_hash)
    groth16_verify(d.vk, redeemer.proof, [hash_scalar])
  }

  else(_) {
    False
  }
}
\end{verbatim}

The datum stores two things: the Poseidon hash of the password (the lock) and
the verification key (the lock's mechanism). Both are public. The redeemer
provides the proof (the key). The key never reveals the password.

\subsection*{Compile}

\begin{verbatim}
aiken check
\end{verbatim}

You should see zero errors and zero warnings.

% ─────────────────────────────────────────────────────────────────────────────
\section{Testing with Aiken} \label{sec:ch6-tests}

Create \texttt{lib/zk\_password/password\_tests.ak} with two tests: one that
asserts a valid proof passes, and one that asserts a wrong hash fails.

\begin{verbatim}
use aiken/crypto/bls12_381/scalar
use zk_password/groth16.{Proof, VerificationKey, groth16_verify}

fn test_vk() -> VerificationKey {
  VerificationKey {
    alpha: #"<vk_alpha1 from compressed_vk.json>",
    beta:  #"<vk_beta2  from compressed_vk.json>",
    gamma: #"<vk_gamma2 from compressed_vk.json>",
    delta: #"<vk_delta2 from compressed_vk.json>",
    ic: [
      #"<vk_ic[0] from compressed_vk.json>",
      #"<vk_ic[1] from compressed_vk.json>",
    ],
  }
}

fn test_proof() -> Proof {
  Proof {
    pi_a: #"<pi_a from compressed_proof.json>",
    pi_b: #"<pi_b from compressed_proof.json>",
    pi_c: #"<pi_c from compressed_proof.json>",
  }
}

// Public signal from public.json (the hash snarkjs used to generate the proof)
const credential_hash: Int =
  9415434746508747181228645764267412889835803815708024285868314342086739058356

test groth16_verify_valid_proof() {
  expect Some(hash_scalar) = scalar.new(credential_hash)
  groth16_verify(test_vk(), test_proof(), [hash_scalar])
}

test groth16_verify_wrong_hash() fail {
  let wrong_hash = credential_hash + 1
  expect Some(hash_scalar) = scalar.new(wrong_hash)
  groth16_verify(test_vk(), test_proof(), [hash_scalar])
}
\end{verbatim}

Fill in the placeholder hex strings with the values from your
\texttt{compressed\_vk.json} and \texttt{compressed\_proof.json}.
The \texttt{credential\_hash} is the integer in \texttt{public.json}.

\begin{verbatim}
aiken check
\end{verbatim}

Expected output:

\begin{verbatim}
Testing ...
  PASS  zk_password/password_tests :: groth16_verify_valid_proof
  PASS  zk_password/password_tests :: groth16_verify_wrong_hash
Summary: 2 passed, 0 failed
\end{verbatim}

\begin{remark}
The \texttt{groth16\_verify\_wrong\_hash} test is marked \texttt{fail}, meaning
it is expected to fail. Aiken marks it as passing precisely \emph{because} the
verification rejects the wrong hash --- confirming the validator cannot be
fooled.
\end{remark}

% ─────────────────────────────────────────────────────────────────────────────
\section{Off-Chain Code with MeshJS} \label{sec:ch6-offchain}

With the validator compiled and tested, we need off-chain code to build the
lock and unlock transactions. We use \textbf{MeshJS}, a TypeScript SDK for
Cardano.

\begin{verbatim}
npm install @meshsdk/core
\end{verbatim}

Read the compiled blueprint and the compressed keys:

\begin{verbatim}
import fs from "fs";
import {
  MeshWallet,
  Transaction,
  resolveScriptHash,
  serializePlutusScript,
  applyParamsToScript,
} from "@meshsdk/core";
import blueprint from "./plutus.json" assert { type: "json" };
import compressedVk from "./compressed_vk.json" assert { type: "json" };
import compressedProof from "./compressed_proof.json" assert { type: "json" };

// Helper: hex string -> Mesh ByteString value
const hex = (s) => ({ alternative: 0, fields: [s] });

// VerificationKey as Plutus Data (mirrors the Aiken record field order)
const vkData = {
  alternative: 0,
  fields: [
    compressedVk.vk_alpha1,
    compressedVk.vk_beta2,
    compressedVk.vk_gamma2,
    compressedVk.vk_delta2,
    compressedVk.vk_ic,   // List<ByteArray>
  ],
};
\end{verbatim}

\subsection*{Locking a UTxO}

To lock funds, build a transaction that sends ADA to the script address with
the datum encoding \texttt{(credential\_hash, vk)}:

\begin{verbatim}
import { publicSignals } from "./public.json" assert { type: "json" };

const credentialHash = publicSignals[0]; // decimal string from public.json

// Datum: Constr(0, [Int(credentialHash), VerificationKey])
const datum = {
  alternative: 0,
  fields: [
    parseInt(credentialHash), // credential_hash as Int
    vkData,                   // vk: VerificationKey
  ],
};

const script = {
  code: blueprint.validators[0].compiledCode,
  version: "V3",
};

const scriptAddress = serializePlutusScript(script, undefined, 0);

const wallet = new MeshWallet({
  networkId: 0,  // 0 = testnet, 1 = mainnet
  fetcher: provider,
  submitter: provider,
  key: { type: "mnemonic", words: yourMnemonic },
});

const tx = new Transaction({ initiator: wallet })
  .sendLovelace(
    { address: scriptAddress, datum: { value: datum, inline: true } },
    "5000000"  // 5 ADA
  );

const signed = await wallet.signTx(await tx.build());
const txHash = await wallet.submitTx(signed);
console.log("Locked:", txHash);
\end{verbatim}

\begin{highlight}
The datum is stored \textbf{inline} on-chain. Anyone can read the Poseidon
hash and the verification key from the UTxO. Neither reveals the password.
\end{highlight}

\subsection*{Unlocking a UTxO}

To unlock, the spender must have generated a valid proof off-chain
(using \texttt{gen\_input.js} + \texttt{snarkjs groth16 fullprove} +
\texttt{compress\_proof.js}). They then build a transaction with the redeemer
encoding the compressed proof:

\begin{verbatim}
// Redeemer: Constr(0, [Proof])
const redeemer = {
  data: {
    alternative: 0,
    fields: [{
      alternative: 0,
      fields: [
        compressedProof.pi_a,
        compressedProof.pi_b,
        compressedProof.pi_c,
      ],
    }],
  },
  budget: { mem: 14000000, cpu: 10000000000 },
};

// Fetch the locked UTxO at the script address
const utxos = await provider.fetchAddressUTxOs(scriptAddress);
const lockedUtxo = utxos[0];

const tx = new Transaction({ initiator: wallet })
  .redeemValue({
    value: lockedUtxo,
    script,
    redeemer,
  })
  .sendLovelace(await wallet.getChangeAddress(), "4800000");

const signed = await wallet.signTx(await tx.build(), true);
const txHash = await wallet.submitTx(signed);
console.log("Unlocked:", txHash);
\end{verbatim}

When this transaction is submitted, the Cardano node evaluates the
\texttt{password\_lock} validator. The validator decompresses the proof and VK,
calls \texttt{groth16\_verify}, and approves the spend only if the pairing
equation holds --- which it will if and only if the spender knows the correct
password.

\begin{highlight}
The password is used only \emph{locally} to run \texttt{fullprove}. It never
appears in the transaction, the redeemer, or anywhere on-chain. An observer who
inspects the blockchain sees the compressed proof and the hash --- neither
reveals the password.
\end{highlight}

% ─────────────────────────────────────────────────────────────────────────────
\section{End-to-End Summary} \label{sec:ch6-summary}

Here is the complete workflow from password to unlocked UTxO:

\begin{enumerate}
  \item \textbf{Setup} (once per circuit): compile \texttt{circuit.circom},
        run the trusted setup ceremony, export \texttt{verification\_key.json},
        compress it to \texttt{compressed\_vk.json} with \texttt{compress\_vk.js}.
  \item \textbf{Lock}: choose a password, run
        \texttt{node circuit/gen\_input.js <password>} to produce
        \texttt{circuit/input.json}, then lock a UTxO with the datum
        \texttt{(credential\_hash, vk)}.
  \item \textbf{Prove}: run \texttt{snarkjs groth16 fullprove} with the
        \texttt{input.json}, the WASM, and \texttt{zkey\_final.zkey}
        to produce \texttt{proof.json}.
  \item \textbf{Compress}: run \texttt{compress\_proof.js} to get
        \texttt{compressed\_proof.json}.
  \item \textbf{Unlock}: build a transaction spending the UTxO with the
        redeemer encoding the compressed proof.
\end{enumerate}

% ─────────────────────────────────────────────────────────────────────────────
\section{How This Differs from Janus Wallet} \label{sec:ch6-vs-janus}

Janus Wallet is a production ZK wallet with staking, delegation, and a
server-side proving API. This chapter's contract is intentionally stripped down
to isolate the core idea.

\begin{center}
\begin{tabular}{lll}
  \hline
  \textbf{Feature}            & \textbf{Janus Wallet}            & \textbf{zk-password} \\
  \hline
  Public inputs               & userId, hash, challenge, flags   & hash only            \\
  Replay protection           & On-chain challenge nonce         & None (educational)   \\
  Server                      & Node.js + PostgreSQL             & Off-chain only       \\
  Hash function               & Poseidon (Circom)                & Poseidon (Circom)    \\
  Curve                       & BLS12-381                        & BLS12-381            \\
  On-chain verifier           & Aiken + builtins                 & Aiken + builtins     \\
  Staking / governance        & Yes                              & No                   \\
  \hline
\end{tabular}
\end{center}

\begin{remark}
The \texttt{zk-password} contract does not include replay protection. In
production, a malicious observer could copy a valid proof from the blockchain
and resubmit it in their own transaction before the legitimate spender. To
prevent this, add a nonce or challenge to the circuit's public inputs and
store an on-chain counter in the datum, as Janus Wallet does with its
\texttt{challenge} and \texttt{challengeFlag} signals.
\end{remark}

This is the complete minimal ZK contract. The next chapters will add replay
protection, multi-input circuits, and a full off-chain proving workflow.
